\title{Controlling Text-to-Image Diffusion by Orthogonal Finetuning}

\begin{abstract}
State-of-the-art text-to-image diffusion models are highly effective at synthesizing realistic images from textual descriptions. A central open challenge is how to steer or constrain these models for various downstream objectives. To this end, we propose a systematic finetuning approach—Orthogonal Finetuning~(OFT)—for adapting text-to-image diffusion architectures to diverse downstream requirements. In contrast to prior finetuning strategies, OFT is theoretically guaranteed to conserve hyperspherical energy, which encodes the relationships among neuron activations projected on the unit hypersphere. Our analysis reveals that maintaining this property is vital to sustaining the semantic generative capacity of text-to-image diffusion models. For enhanced training robustness, we introduce Constrained Orthogonal Finetuning (COFT), which integrates an explicit radial constraint onto the hypersphere. Our investigation concentrates on two pivotal text-to-image adaptation problems: subject-driven generation—where, provided with a small sample of subject images and a textual prompt, the aim is to produce new images of the same subject in novel settings—and controllable generation, wherein the diffusion model is directed by external control inputs. Empirical studies demonstrate that the OFT methodology surpasses existing alternatives in both image generation quality and training efficiency.
\end{abstract}

\section{Introduction}

Text-to-image diffusion models~\cite{saharia2022photorealistic,ramesh2022hierarchical,rombach2022high} have recently set new benchmarks for high-resolution, faithful image synthesis conditioned on textual prompts. However, despite these advances, guiding image generation with text alone often leads to ambiguity and does not suffice for precise or detailed manipulations of output images. To overcome this limitation, we address two central text-to-image generation paradigms in this work:

\begin{itemize}[leftmargin=*,nosep]

    \item \textbf{Subject-driven generation}~\cite{ruiz2023dreambooth}: The challenge here is, given a handful of images featuring a specific subject, to synthesize new depictions of the same subject placed in varied scenes or contexts using textual guidance.
    \item \textbf{Controllable generation}~\cite{zhang2023adding,mou2023t2i}: In this setting, the task involves producing images that adhere to auxiliary control inputs (\emph{e.g.}, edge maps, segmentation masks) alongside text prompts.
\end{itemize}

At their core, both tasks focus on the challenge of finetuning text-to-image diffusion models while maintaining the generative quality achieved during pretraining. We outline two main requirements for an optimal finetuning technique: (1) \emph{training efficiency}, characterized by minimizing both the count of trainable parameters and the number of training epochs, and (2) \emph{generalizability preservation}, ensuring that the model continues to produce high-fidelity and diverse outputs. Common approaches to finetuning involve either making slight adjustments to the neuron weights using a low learning rate (\emph{e.g.}, \cite{ruiz2023dreambooth}) or introducing a lightweight module that re-parameterizes certain weights (\emph{e.g.}, \cite{hulora2022,zhang2023adding}). Despite their straightforward nature, these strategies do not inherently safeguard the model’s generative capabilities learned during pretraining. Furthermore, there is a noticeable gap in systematic methodologies for developing effective finetuning protocols and identifying appropriate hyperparameter values, such as the optimal number of epochs. One significant challenge is the absence of a robust metric to assess the retention of the model's pretrained generative abilities. Many current finetuning practices operate under the implicit premise that minimizing the Euclidean distance between finetuned and pretrained model weights correlates with better preservation of prior performance, which leads to the routine use of very small learning rates. However, we contend that Euclidean proximity to the pretrained parameters alone is insufficient to capture true semantic preservation. Accordingly, leveraging a more structurally informed metric to compare the finetuned and pretrained models could greatly enhance both the consistency of finetuning and the maintenance of generative performance.

Drawing from empirical findings indicating that hyperspherical similarity effectively captures semantic relationships~\cite{liu2017deep,liu2018decoupled,chen2020angular}, our approach leverages hyperspherical energy~\cite{liu2018learning} to quantify the pairwise interactions between neurons within a layer. Here, hyperspherical energy is calculated as the aggregate of hyperspherical similarities (\emph{e.g.}, cosine similarity) among every neuron pair in a given layer, thus reflecting the distribution uniformity of neurons over the unit hypersphere~\cite{liu2021learning}. We posit that an optimally finetuned model should exhibit minimal deviation in hyperspherical energy from its pretrained counterpart. One straightforward approach would be to introduce a regularization term that maintains the hyperspherical energy unchanged during finetuning; however, this strategy does not necessarily ensure a minimized hyperspherical energy difference. To address this, we exploit the invariance property inherent to hyperspherical energy: specifically, when all neurons undergo an identical orthogonal transformation, their pairwise hyperspherical similarity remains unchanged by design. This observation underpins our proposal of Orthogonal Finetuning~(OFT), a technique for adapting large-scale text-to-image diffusion models to specific tasks without altering their hyperspherical energy. The fundamental premise of OFT is to learn an orthogonal transformation that is shared across the entire layer, thereby preserving the angular relationships among neurons. In essence, OFT can be interpreted as transforming the coordinate basis for the neurons within a layer. Further, recognizing that reduced Euclidean distance between the finetuned and original model indicates stronger retention of pretrained capabilities, we introduce a refined variant, Constrained Orthogonal Finetuning~(COFT), which additionally restricts the finetuned solution to remain within a hyperspherical shell of fixed radius around the pretrained neuron representations.

The rationale behind employing orthogonal transformation for neuron finetuning is partially inspired by the 2D Fourier transform, where an image can be represented by its magnitude and phase spectra. The phase spectrum—encoding the angular relation between the input and the basis—retains the essential semantic content. As an illustration, reconstructing an image using its original phase spectrum paired with an arbitrary magnitude spectrum still maintains the core semantics of the image. This observation implies that altering the directions of neurons fundamentally enables semantic adaptation of the generated image, which underpins both subject-driven and controllable generation objectives. Nevertheless, allowing unrestricted modification of neuron orientations risks impairing the pretrained model's generative capability. To address this, we advocate preserving the relative angles among all neuron pairs, predicated on the hypothesis that such angular relationships are fundamental for encoding neural network knowledge. Consequently, we learn a shared orthogonal transformation per layer to maintain invariant hyperspherical energy, aligning with our geometric intuition.

Additionally, our approach is influenced by orthogonal over-parameterized training~\cite{liu2021orthogonal}, wherein classification models are trained from scratch by applying orthogonal transformations to a randomly initialized network. A randomly initialized neural network exhibits a predictably low hyperspherical energy on average, and \cite{liu2021orthogonal} seeks to maintain this low energy throughout training, since reduced energy correlates with enhanced generalization in classification~\cite{liu2018learning,lin2020regularizing}. Their findings indicate that orthogonal transformations offer sufficient expressiveness to train broadly generalizable classification models. Distinctly, our focus lies in adapting text-to-image diffusion models for greater controllability and improved generative performance on downstream tasks. The distinction between OFT and \cite{liu2021orthogonal} is twofold: Firstly, \cite{liu2021orthogonal} aims to further reduce hyperspherical energy, whereas OFT seeks to conserve the hyperspherical energy present in the pretrained model, thereby safeguarding intrinsic learned structure during finetuning. Direct minimization of hyperspherical energy in diffusion models could undermine pre-existing semantic organization. Secondly, OFT prioritizes minimizing deviation from the pretrained parameters, naturally yielding a constrained form; conversely, \cite{liu2021orthogonal} operates without such constraints. The central challenge in finetuning is to balance adaptability with stability, and we posit that OFT effectively navigates this trade-off. Below, we outline our primary contributions:

\begin{itemize}[leftmargin=*,nosep]

    \item We introduce a new finetuning technique, Orthogonal Finetuning, designed to enhance the controllability of text-to-image diffusion models. To further ensure stable optimization, we develop a constrained version that restricts the angular change relative to the pretrained model.
    \item Unlike conventional finetuning strategies, OFT enables model adaptation while formally guaranteeing the preservation of hyperspherical energy—a property we observe to be crucial for retaining the generative semantics of the original pretrained model.
    \item We deploy OFT across two settings: subject-driven generation and controllable generation. Through extensive experiments, we show that our method substantially surpasses existing approaches in generation fidelity, faster convergence, and increased robustness during finetuning. In addition, OFT demonstrates superior sample efficiency by reaching satisfactory performance with far fewer training samples and epochs.
    \item For the task of controllable generation, we propose a novel control consistency metric to quantitatively assess controllability. This metric operates by extracting the control signal from the model’s generated output and measuring its similarity to the original control input. Our results confirm that OFT achieves notably stronger controllability according to this criterion.
\end{itemize}

\section{Related Work}

\textbf{Text-to-image diffusion models}. The field of text-to-image synthesis has rapidly advanced~\cite{saharia2022photorealistic,ramesh2022hierarchical,rombach2022high,gu2022vector,nichol2022glide}, primarily due to significant breakthroughs in diffusion models~\cite{ho2020denoising,song2021score,song2021denoising,dhariwal2021diffusion} and cross-modal vision-language learning~\cite{su2020vlbert,lu2019vilbert,radford2021learning,wang2021simvlm,li2022blip,li2023blip,alayrac2022flamingo,schuhmann2022laion}. For example, GLIDE~\cite{nichol2022glide} and Imagen~\cite{saharia2022photorealistic} operate in the image pixel domain; GLIDE employs joint training of the text encoder and a diffusion prior, leveraging text-image pairs, while Imagen relies on a frozen text encoder pretrained elsewhere. In contrast, Stable Diffusion~\cite{rombach2022high} and DALL-E2~\cite{ramesh2022hierarchical} implement diffusion processes in the latent space. Stable Diffusion utilizes VQ-GAN~\cite{esser2021taming} to derive a compact visual representation, whereas DALL-E2 integrates CLIP~\cite{radford2021learning} to define a shared embedding space for both modalities. Beyond diffusion models, generative adversarial networks~\cite{reed2016generative,zhang2017stackgan,xu2018attngan,li2019controllable} and autoregressive models~\cite{ramesh2021zero,ding2021cogview,wu2022nuwa,yuscaling} are also explored for text-to-image tasks. OFT is inherently model-agnostic and applicable to a broad class of text-to-image diffusion architectures.

\textbf{Subject-driven generation}. Approaches such as~\cite{avrahami2022blended,nichol2022glide} utilize user-provided masks to prevent alteration of target subjects during image generation. Techniques based on inversion~\cite{choi2021ilvr,dhariwal2021diffusion,ramesh2022hierarchical,gal2022image} enable modifying contextual elements while maintaining the original subject intact. The method in~\cite{hertz2022prompt} allows for both localized and global edits without the need for external masks. However, these methods often struggle to preserve intricate, identity-specific characteristics. Pivotal Tuning~\cite{roich2022pivotal} adapts a generator near an initial inverted latent representation using regularization to maintain subject identity, while~\cite{nitzan2022mystyle} learns personalized facial priors from a collection of images per individual. Casanova et al.~\cite{casanova2021instance} can generate multiple variations of an instance, though possibly at the expense of instance-level details. DreamBooth~\cite{ruiz2023dreambooth} fine-tunes diffusion models using both reconstruction and prior preservation losses, starting from limited subject images and a custom token. While OFT operates within the DreamBooth paradigm, it differentiates itself by replacing standard low-rate finetuning with updates based on orthogonal transformations.

\textbf{Controllable generation}. Image-to-image translation can be interpreted as an instance of controllable generation. Traditional approaches predominantly rely on conditional generative adversarial networks~\cite{isola2017image,zhu2017toward,wang2018high,park2019semantic,choi2018stargan} for this purpose. Diffusion models have also gained traction in image-to-image translation tasks~\cite{saharia2022palette,voynov2022sketch,wang2022pretraining}. Among recent advancements, ControlNet~\cite{zhang2023adding} stands out by enabling a pretrained diffusion model to respond to supplementary control signals through finetuning and adaptation, attaining remarkable results in controllable generation. Similarly, T2I-Adapter~\cite{mou2023t2i} independently introduces stronger control over outputs by finetuning diffusion models. Following the protocol established in \cite{zhang2023adding,mou2023t2i}, we employ OFT to adjust pretrained diffusion models. This approach consistently enhances controllability with reduced requirements for training data and finetuning parameters. Importantly, OFT introduces no extra computational cost at inference time.

\textbf{Model finetuning}. The strategy of refining large pretrained models for downstream applications has become increasingly widespread~\cite{devlin2018bert,brown2020language,he2022masked}. Within this context, adaptation techniques (\emph{e.g.}, \cite{hulora2022,houlsby2019parameter,pfeiffer2020adapterfusion}) are extensively explored in natural language processing. OFT is closely related to LoRA~\cite{hulora2022}, which applies low-rank updates additively during model finetuning. Unlike LoRA, OFT employs a multiplicative, layer-shared orthogonal transformation to modify neuron weights. This method ensures that the angles between neurons within a layer remain unchanged, providing improved stability.

\section{Orthogonal Finetuning}

\subsection{Why Does Orthogonal Transformation Make Sense?}

To motivate the use of orthogonal transformations in the context of finetuning text-to-image diffusion models, we break down the rationale into two key aspects: (1) the motivation for tuning neuron angles (i.e., direction), and (2) the justification for using orthogonal transformations specifically for this purpose.

To address the first question, we build upon prior empirical findings from \cite{liu2018decoupled,chen2020angular} which indicate that angular disparities in features provide a reliable measure of semantic distance. SphereNet~\cite{liu2017deep} demonstrates that constraining all neuron activations to a unit hypersphere does not diminish representational power and can actually enhance generalization performance, suggesting that neuron direction encodes critical data characteristics. To further highlight the significance of neuron angularity, we designed a simple illustrative experiment (see Figure~\ref{fig:toy_ae}), where a basic convolutional autoencoder is trained on images of flowers. During training, we generate the feature map using the standard inner product ($z$ is the convolution result from kernel $\bm{w}$ and window input $\bm{x}$). At test time, we contrast three methods of obtaining the feature map: (a) inner product as during training, (b) utilizing solely the norm, and (c) employing only the angle information. Results in Figure~\ref{fig:toy_ae} indicate that neuron angles alone are nearly sufficient for perfect image reconstruction, while the norms appear to encode no relevant content. Notably, the cosine activation (case c) is not used during the training phase—the network is optimized exclusively through inner products. These outcomes suggest that neuron directions are principally responsible for preserving the semantic content of images; hence, adjusting neuron directions offers a promising avenue for semantic modification through finetuning.

Turning to the second question, the most direct approach to adjusting neuron directions is to jointly rotate or reflect all neurons within a layer, a procedure naturally expressed as an orthogonal transformation. While it is possible to consider more granular angular manipulations by rotating neurons individually, our experiments indicate that orthogonal transformations strike an optimal balance between expressiveness and structural consistency. Furthermore, as shown in \cite{liu2021orthogonal}, orthogonal transformations possess ample expressive capacity for neural network training. To empirically corroborate this, we conduct an experiment—using the ControlNet~\cite{zhang2023adding} protocol—to assess the regularizing impact of orthogonality. Figure~\ref{fig:ortho} reports that our canonical OFT consistently delivers robust and precise control upon completion of training (by epoch 20), whereas an OFT implementation lacking the orthogonality constraint produces neither plausible images nor controllability. This experiment underscores the critical role of enforcing orthogonality in OFT.

\subsection{General Framework}

The standard approach for finetuning generally relies on applying gradient descent with a relatively small learning rate, which serves to update the model parameters—or just specific layers—incrementally. By employing a small learning rate, this procedure effectively encourages the adjusted parameters to remain close to those of the original pretrained model. In essence, typical finetuning attempts to optimize the model while implicitly minimizing $\thickmuskip=2mu \medmuskip=2mu \| \bm{M} - \bm{M}^0\|$, where $\bm{M}$ and $\bm{M}^0$ represent the weight matrices of the finetuned and pretrained models, respectively. While this implicit regularization provides some intuitive constraint, it may permit excessive flexibility, especially in the context of large-scale models. To strengthen this, LoRA incorporates an explicit low-rank limitation on the parameter update, specifically enforcing $\thickmuskip=2mu \medmuskip=2mu \text{rank}(\bm{M} - \bm{M}^0)=r'$ with $r'$ being a modest integer. In contrast, OFT imposes a restriction based on neuron similarity, formalized as $\thickmuskip=2mu \medmuskip=2mu \| \text{HE}(\bm{M})-\text{HE}(\bm{M}^0) \|=0$, where $\text{HE}(\cdot)$ indicates the hyperspherical energy for the matrix. 

To elucidate, consider a fully connected layer denoted as $\thickmuskip=2mu \medmuskip=2mu \bm{W}=\{\bm{w}_1,\cdots,\bm{w}_n\}\in\mathbb{R}^{d\times n}$, where each column vector $\bm{w}_i\in\mathbb{R}^{d}$ corresponds to the $i$-th neuron, and $\bm{W}^0$ is the pretrained layer weight matrix. The layer output $\thickmuskip=2mu \medmuskip=2mu \bm{z}\in\mathbb{R}^n$ is calculated by $\thickmuskip=2mu \medmuskip=2mu \bm{z}=\bm{W}^\top\bm{x}$, with $\thickmuskip=2mu \medmuskip=2mu \bm{x}\in\mathbb{R}^d$ as the input. Under OFT, the goal becomes minimizing the change in hyperspherical energy between the updated and initial parameters:

\begin{equation}\label{eq:oft_principle}
\footnotesize
\min_{\bm{W}} \left\| \text{HE}(\bm{W})-\text{HE}(\bm{W}^0)\right\|~~\Leftrightarrow~~\min_{\bm{W}} \bigg{\|} \sum_{i\neq j} \|\hat{\bm{w}}_i-\hat{\bm{w}}_j\|^{-1}-\sum_{i\neq j} \|\hat{\bm{w}}^0_i-\hat{\bm{w}}^0_j\|^{-1}\bigg{\|}
\end{equation}

where the normalized neurons are defined by $\thickmuskip=2mu \medmuskip=2mu \hat{\bm{w}}_i:=\bm{w}_i/\|\bm{w}_i\|$, and $\text{HE}(\bm{W}):= \sum_{i\neq j} \|\hat{\bm{w}}_i-\hat{\bm{w}}_j\|^{-1}$ specifies the hyperspherical energy for $\bm{W}$. Noticeably, Eq.~\eqref{eq:oft_principle} attains its minimum value of zero when $\bm{W}$ and $\bm{W}^0$ are connected by either a rotation or reflection—that is, if $\thickmuskip=2mu \medmuskip=2mu \bm{W}=\bm{R}\bm{W}^0$ with $\thickmuskip=2mu \medmuskip=2mu \bm{R}\in\mathbb{R}^{d\times d}$ being orthogonal (with determinant $1$ or $-1$, corresponding to rotation or reflection, respectively). The central concept in OFT, therefore, is to restrict model updates to learning orthogonal matrices that transform the neuron representations in each layer. Formally, for a pretrained fully connected layer $\thickmusk

\begin{equation}\label{eq:oft_general}
    \footnotesize
    \bm{z}=\bm{W}^\top\bm{x}=(\bm{R}\cdot\bm{W}^0)^\top\bm{x},~~~\text{s.t.}~\bm{R}^\top\bm{R}=\bm{R}\bm{R}^\top=\bm{I}
\end{equation}

In this expression, $\bm{W}$ refers to the weight matrix updated via OFT, and $\bm{I}$ represents the identity matrix. Figure~\ref{fig:block} provides a visual depiction of OFT. As in LoRA’s approach with zero initialization, it is important for OFT to adaptively finetune the model, starting exactly from $\bm{W}^0$. To realize this, the orthogonal matrix $\bm{R}$ is initially set to the identity, which ensures the finetuned network coincides with the pretrained parameters at the outset. To maintain the orthogonality constraint on $\bm{R}$ throughout training, differentiable orthogonalization methods from \cite{liu2021orthogonal,lezcano2019cheap} can be applied. Computationally efficient strategies to enforce this property will be examined below.

\subsection{Efficient Orthogonal Parameterization}\label{sect:eff_ortho}

Traditional orthogonalization approaches like the differentiable Gram-Schmidt process~\cite{liu2021orthogonal} can incur prohibitive computational costs for large models. To improve efficiency, we opt for the Cayley parameterization when constructing the orthogonal matrix. Specifically, $\bm{R}$ is defined as $\thickmuskip=2mu \medmuskip=2mu \bm{R}=(\bm{I}+\bm{Q})(I-\bm{Q})^{-1}$, with $\bm{Q}$ being a skew-symmetric matrix, i.e., $\thickmuskip=2mu \medmuskip=2mu \bm{Q}=-\bm{Q}^\top$. While this Cayley parameterization restricts $\bm{R}$ to orthogonal matrices with determinant $1$ (members of the special orthogonal group), in our empirical analysis this restriction poses no observable limitations. However, even under the Cayley parameterization, the number of parameters in $\bm{R}$ grows rapidly with larger $d$, making the approach parameter-inefficient for high-dimensional matrices. To mitigate this, we adopt a block-diagonal structure for $\bm{R}$, decomposing it into $r$ orthogonal blocks, resulting in the following structure:

\begin{equation}
    \footnotesize
    \bm{R}=\text{diag}(\bm{R}_1,\bm{R}_2,\cdots,\bm{R}_r)=\begin{bmatrix}
    \bm{R}_1\in \text{O}(\frac{d}{r}) &  &\\
    &\ddots & \\
    & & \bm{R}_r\in \text{O}(\frac{d}{r})
    \end{bmatrix}\in \text{O}(d)
\end{equation}

where $\text{O}(d)$ refers to the set of $d$-dimensional orthogonal matrices, and $\thickmuskip=2mu \medmuskip=2mu \bm{R}\in\mathbb{R}^{d\times d}$ as well as $\thickmuskip=2mu \medmuskip=2mu \bm{R}_i\in\mathbb{R}^{d/r\times d/r}$ for each $i$ are orthogonal matrices. When $\thickmuskip=2mu \medmuskip=2mu r=1$, the block-diagonal structure simplifies to a fully unconstrained orthogonal matrix. For a general $d\times d$ orthogonal matrix, the number of free parameters is $\thickmuskip=2mu \medmuskip=2mu d(d-1)/2$, which corresponds to $\mathcal{O}(d^2)$ complexity. In comparison, for an $r$-block diagonal orthogonal matrix, the parameter count becomes $\thickmuskip=2mu \medmuskip=2mu d(d/r-1)/2$, scaling as $\thickmuskip=2mu \medmuskip=2mu \mathcal{O}(d^2/r)$. Additionally, if all block matrices are shared—i.e., $\thickmuskip=2mu \medmuskip=2mu\bm{R}_i=\bm{R}_j$ for all $i\neq j$—then the total parameter complexity drops to $\mathcal{O}(d^2/r^2)$. These techniques allow parameter count reduction without compromising the orthogonality of $\bm{R}$, so the hyperspherical energy is still preserved.

We next consider the parameter efficiency of OFT relative to LoRA. For LoRA, given a low-rank setting with rank $r'$, the number of tunable parameters is $\thickmuskip=2mu \medmuskip=2mu r'(d+n)$. Assuming both $r$ and $r'$ are proportional to $d$, such that $\thickmuskip=2mu \medmuskip=2mu r = r' = \alpha d$ for some constant $\thickmuskip=2mu \medmuskip=2mu 0 < \alpha \leq 1$, the LoRA parameter count becomes $\mathcal{O}(d^2 + dn)$, while the complexity for OFT reduces to $\mathcal{O}(d)$. To highlight these differences, consider a weight matrix with dimensions $\thickmuskip=2mu \medmuskip=2mu 128\times128$. In this case, LoRA with $\thickmuskip=2mu \medmuskip=2mu r' = 8$ yields $2,048$ parameters, while OFT with $\thickmuskip=2mu \medmuskip=2mu r = 8$ (without block sharing) only requires $960$ parameters.

\subsection{Constrained Orthogonal Finetuning}

The expressiveness of standard OFT can be further limited by constraining the finetuned parameters to stay close to those of the pretrained model. COFT realizes this restriction with the following formulation for the forward computation:

\begin{equation}\label{eq:coft_general}
    \footnotesize
    \bm{z}=\bm{W}^\top\bm{x}=(\bm{R}\cdot\bm{W}^0)^\top\bm{x},~~~\text{s.t.}~\bm{R}^\top\bm{R}=\bm{R}\bm{R}^\top=\bm{I},~\norm{\bm{R}-\bm{I}}\leq \epsilon
\end{equation}

where both an orthogonality condition and a bound $\epsilon$ on the deviation from the identity matrix are enforced. Orthogonality can be maintained using the Cayley transform, as discussed in Section~\ref{sect:eff_ortho}. Nonetheless, implementing the additional $\epsilon$-deviation constraint within the Cayley parameterization poses challenges. To further analyze the behavior of the Cayley transform, we employ the Neumann series for the approximation $\thickmuskip=2mu \medmuskip=2mu \bm{R}=(\bm{I}+\bm{Q})(I-\bm{Q})^{-1} \approx \bm{I} + 2\bm{Q} + \mathcal{O}(\bm{Q}^2)$ (assuming the Ne

Since the series converges under the operator norm, it is valid to transfer the constraint $\thickmuskip=2mu \medmuskip=2mu\|\bm{R}-\bm{I}\|\leq\epsilon$ directly inside the Cayley transform. This reformulates the constraint as $\thickmuskip=2mu \medmuskip=2mu\|\bm{Q}-\bm{0}\|\leq\epsilon'$, where $\bm{0}$ represents a zero matrix and $\epsilon'$ is a distinct tolerance parameter from $\epsilon$. Projected gradient descent allows us to enforce the new constraint on $\bm{Q}$ with ease. To initialize the orthogonal matrix $\bm{R}$ to the identity, we begin with $\bm{Q}$ set to the zero matrix. The COFT methodology, therefore, incorporates two explicit forms of regularization: ensuring a minimal hyperspherical energy gap and maintaining a controlled deviation from the initial pretrained parameters. Whereas the latter constraint is often implicitly assumed in current finetuning approaches, COFT treats it as an explicit constraint. Empirically, we find that although OFT already performs strongly, COFT introduces greater stability during the finetuning phase, which can be attributed to the explicit deviation constraint. Figure~\ref{fig:eps_coft} demonstrates the influence of $\epsilon$ on COFT's effectiveness. As seen in the figure, $\epsilon$ acts as a knob that regulates the degree of adaptability during finetuning: increasing $\epsilon$ yields behavior more akin to OFT, while decreasing $\epsilon$ brings the COFT-tuned model closer to the pretrained text-to-image diffusion model in function.

\subsection{Re-scaled Orthogonal Finetuning}

We present an enhancement to standard OFT by incorporating a trainable scaling parameter for each neuron. The rationale is that scaling neuron activations leaves the hyperspherical energy unchanged, as normalization removes magnitude effects. The forward computation in our modification is: $\thickmuskip=2mu \medmuskip=2mu\bm{z}=(\bm{R}\bm{W}^0\bm{D})^\top\bm{x}$,\footnote{\textbf{Errata}: The NeurIPS camera-ready paper incorrectly states this as $\thickmuskip=2mu \medmuskip=2mu\bm{z}=(\bm{D}\bm{R}\bm{W}^0)^\top\bm{x}$. The initial implementation, however, uses the correct formulation, so Appendix~\ref{app:rescaled}'s findings remain valid.} where $\thickmuskip=2mu \medmuskip=2mu \bm{D} = \text{diag}(s_1,\cdots,s_n) \in \mathbb{R}^{n\times n}$ is a diagonal matrix whose entries $s_1,\ldots,s_n$ are all strictly positive and subject to learning. Unlike the original OFT, which only tunes $\bm{R}$ in Eq.~\eqref{eq:oft_general}, this variant learns both the diagonal scaling matrix $\bm{D}$ and the orthogonal matrix $\bm{R}$. The proposed re-scaled OFT thus increases OFT's expressive capacity with minimal overhead in the number of parameters. For experimental clarity, we retain the baseline OFT to focus solely on orthogonal transformations, but our evaluations indicate that the re-scaled variant often provides further gains (see Appendix~\ref{app:rescaled}).

\section{Intriguing Insights and Discussions}

\textbf{OFT is compatible with various neural network designs}. OFT is, by its formulation, applicable to any neural architecture. In the context of Transformers, common practice with LoRA is to target the attention matrix weights~\cite{hulora2022}. To ensure an equitable comparison with LoRA, our experiments restrict OFT's use to attention weight finetuning as well. Beyond linear layers, OFT's design makes it particularly suitable for convolutional layers. Here, the block-diagonal format of $\bm{R}$ yields distinct advantages not seen with LoRA. If we select the number of blocks to match the input channel count, each block governs the transformation for an individual channel, paralleling the mechanics of depth-wise convolution kernels~\cite{chollet2017xception}. If $\bm{R}$'s blocks are shared, the same orthogonal transformation is applied across all channels. Further, preliminary experimental results regarding OFT on convolutional layers are included in Appendix~\ref{app:convolution}.

\textbf{Connection to LoRA}. LoRA introduces a low-rank matrix to avoid major alterations of the pretrained weight matrix, thus preserving the learned representations. In comparison, OFT enforces the transform applied to the pretrained weight matrix to be orthogonal (and full-rank), which safeguards the integrity of information acquired during pretraining by constraining transformations that could degrade it. The OFT forward computation can be reformulated as $\thickmuskip=2mu \medmuskip=2mu \bm{z}=(\bm{R}\bm{W}^0)^\top\bm{x}=(\bm{W}^0+(\bm{R}-\bm{I})\bm{W}^0)^\top\bm{x}$, where the term $\thickmuskip=2mu \medmuskip=2mu (\bm{R}-\bm{I})\bm{W}^0$ serves a role akin to the low-rank update in LoRA. Since $\bm{W}^0$ usually has full rank, OFT achieves a low-rank update as well when $\thickmuskip=2mu \medmuskip=2mu \bm{R}-\bm{I}$ is low-rank. Much like LoRA, which utilizes a rank hyperparameter $r'$, OFT makes use of a diagonal block parameter $r$ to control trainable parameter count. Notably, LoRA and OFT provide two orthogonal approaches to parameter efficiency: LoRA capitalizes on low-rank structures to cut down the number of learned parameters, whereas OFT achieves efficiency through leveraging sparsity, specifically by utilizing block-diagonal orthogonal matrices.

\textbf{Why OFT converges faster}. As illustrated in Figure~\ref{fig:toy_ae}, changing the orientations of neurons most effectively alters the encoded semantic content—a process that OFT is well-suited to enable. Additionally, OFT's methodology can be interpreted as the finetuning of neurons constrained to a smooth hypersphere manifold, which contributes to a more favorable optimization landscape. Empirical observations supporting this are discussed in \cite{liu2021orthogonal}.

\textbf{Why not minimize hyperspherical energy}. A significant distinction from \cite{liu2021orthogonal} is that we do not pursue hyperspherical energy minimization. In classification settings, it is important that neurons are non-redundant; achieving minimum hyperspherical energy results in neurons being evenly distributed across the hypersphere, which is not necessarily a productive finetuning objective since it risks obliterating valuable pretraining knowledge.

\textbf{Balancing flexibility and regularization in finetuning}. Our analysis reveals a fundamental tension between flexibility and regularization in finetuning strategies. Conventional finetuning grants maximal flexibility, yet often suffers from instability and is prone to model collapse. OFT, though conceptually straightforward, achieves an effective compromise between adaptability and stability by maintaining the relative angles between neuron pairs. Furthermore, using a block-diagonal structure to parametrize the orthogonal matrix can be interpreted as enforcing even stronger regularization.

\textbf{No increase in inference cost}. In contrast to methods such as ControlNet, our OFT approach incurs no additional computational expense during inference. At test time, the model simply incorporates the learned orthogonal matrix $\bm{R}$ with the pretrained weight matrix $\bm{W}^0$ to form an updated weight matrix $\thickmuskip=2mu \medmuskip=2mu \bm{W}=\bm{R}\bm{W}^0$. This procedure ensures that the inference throughput remains unchanged compared to the original pretrained architecture.

\section{Experiments and Results}

\textbf{Experimental setup}. Our experiments utilize Stable Diffusion v1.5~\cite{rombach2022high} as the foundational text-to-image model. For equitable comparison, generated samples are randomly selected across all evaluated methods. In subject-driven image synthesis, our protocol aligns with the procedures of DreamBooth~\cite{ruiz2023dreambooth}. For tasks requiring explicit control, we adopt experimental settings inspired by ControlNet~\cite{zhang2023adding} and T2I-Adapter~\cite{mou2023t2i}. To directly compare with LoRA, we ensure OFT or COFT modifications target the same layers as LoRA. Further implementation details and supplementary results can be found in Appendix~\ref{app:settings}.

\subsection{Subject-driven Generation}

\textbf{Protocol}. DreamBooth~\cite{ruiz2023dreambooth} and LoRA~\cite{hulora2022} serve as primary baselines in this evaluation. All approaches use the DreamBooth loss function for training. For both DreamBooth and LoRA, we closely adhere to the respective original methodologies and optimize hyperparameters as recommended. Additional empirical results are available in Appendices~\ref{app:settings},\ref{app:coft_oft},\ref{app:more_qualitative},\ref{app:failure}.

\textbf{Finetuning stability and convergence}. We begin by assessing the stability and convergence rate of finetuning across DreamBooth, LoRA, OFT, and COFT. Figures~\ref{fig:teaser} and~\ref{fig:db_convergence} illustrate the comparative results. Notably, both COFT and OFT maintain stable finetuning behavior when applied to the diffusion model. By the 400th iteration, DreamBooth and OFT-based approaches demonstrate effective control over the generated content, whereas LoRA is unable to maintain subject identity. After reaching 2000 iterations, DreamBooth starts to exhibit mode collapse in its outputs and LoRA misrepresents the attribute by rendering yellow fur rather than a yellow shirt. On the other hand, OFT and COFT continue to provide stable and reliable manipulation over image generation at this stage. These outcomes reinforce the notion that our OFT framework achieves both rapid and robust convergence in subject-driven generation tasks. The minimal sensitivity to the choice of finetuning iterations further highlights the practical advantages of OFT and COFT, as it reduces the need to fine-tune the number of training iterations for different subjects. Thus, it is feasible to set a relatively high iteration count for OFT and COFT without requiring careful adjustment. In the case of COFT, with an appropriately chosen $\epsilon$, tuning both the learning rate and the iteration number becomes a straightforward task.

\textbf{Quantitative comparison}. Adopting the evaluation protocols from \cite{ruiz2023dreambooth}, we quantitatively measure subject fidelity (using DINO~\cite{caron2021emerging}, CLIP-I~\cite{radford2021learning}), text prompt fidelity (CLIP-T~\cite{radford2021learning}), and output diversity (LPIPS~\cite{zhang2018unreasonable}). The CLIP-I metric is based on the mean pairwise cosine similarity between CLIP embeddings of generated and actual images. DINO follows the same procedure, substituting ViT S/16 DINO embeddings. CLIP-T calculates the mean cosine similarity between CLIP embeddings of the textual prompt and those of the generated images. Additionally, the average LPIPS cosine similarity between images of the same subject generated from the same prompt serves as our measure of diversity. As presented in Table~\ref{table:dreambooth_final}, COFT and OFT significantly surpass DreamBooth and LoRA in the DINO and CLIP-I evaluations, and display equal or improved performance in prompt fidelity and diversity. Each algorithm undergoes finetuning with the same pretrained model across 30 distinct random seeds to mitigate the influence of randomness. These findings indicate that our OFT approach not only ensures more stable and faster convergence, but also consistently delivers superior performance in the final results.

\textbf{Qualitative comparison}. To gain a clearer and more direct sense of the advantages brought by OFT, we present several randomly selected cases of subject-driven image synthesis in Figure~\ref{fig:dreambooth_final}. For an equitable evaluation, all results are produced by applying each approach to the same finetuned model, with no custom prompt engineering for individual images. For each compared method, we use the checkpoint with the highest validation CLIP scores. Figure~\ref{fig:dreambooth_final} reveals that both OFT and COFT effectively maintain the semantic features of the subject, whereas LoRA frequently struggles to retain subject identity (\emph{e.g.}, LoRA entirely fails in preserving subject appearance in the bowl instance). Additionally, OFT and COFT display substantially stronger prompt adherence, whereas DreamBooth, despite keeping the subject intact, often generates images that do not align with the prompt content (\emph{e.g.}, see the first row of the bowl case). These qualitative observations underscore OFT’s simultaneous advantages in prompt responsiveness and subject fidelity. Furthermore, OFT’s performance is robust to the number of training iterations, consistently delivering strong results even with extensive finetuning—a property not observed for DreamBooth or LoRA. Additional qualitative illustrations are provided in Appendix~\ref{app:more_qualitative}. In addition, Appendix~\ref{app:human} presents a human evaluation corroborating OFT’s superior performance.

\subsection{Controllable Generation}

\textbf{Settings}. For baseline comparison, we utilize ControlNet~\cite{zhang2023adding}, T2I-Adapter~\cite{mou2023t2i}, and LoRA~\cite{hulora2022}. The primary evaluation encompasses three demanding tasks of controllable image generation: Canny edge-to-image~(C2I) conversion on COCO~\cite{lin2014microsoft}, segmentation-to-image~(S2I) on ADE20K~\cite{zhou2017scene}, and landmark-to-face~(L2F) synthesis using CelebA-HQ~\cite{karras2018progressive,xia2021tedigan}. Each method is finetuned from Stable Diffusion~(SD) v1.5 for 20 training epochs on the respective datasets. Expanded results appear in Appendix~\ref{app:more_qualitative},\ref{app:more_control_task},\ref{app:failure}.

\textbf{Convergence}. We assess the rate of convergence for ControlNet, T2I-Adapter, LoRA, and COFT on the L2F benchmark through both quantitative and qualitative means. For our metric, we calculate the average $\ell_2$ distance between the predicted face landmarks and those specified by the control input. In Figure~\ref{fig:face_convergence}, we depict how face landmark error varies as each model is finetuned over a range of training epochs. It is evident that COFT and OFT attain faster convergence than the competing approaches. Specifically, LoRA requires 20 epochs to match the performance that OFT achieves as early as the 8-th epoch. It should be noted that both OFT and COFT are parameter-efficient, utilizing a number of trainable parameters comparable to LoRA and substantially fewer than ControlNet, yet they still reach optimal performance more rapidly than previous methods. Additionally, Figure~\ref{fig:teaser} further substantiates the rapid convergence of OFT: the right panel demonstrates that OFT significantly outperforms ControlNet and LoRA in terms of data efficiency, attaining strong convergence with only 5\% of the ADE20K dataset. For qualitative comparisons, our analysis centers on OFT, COFT, and ControlNet—since ControlNet most closely matches our landmark error. Figure~\ref{fig:face_convergence_qual} reveals that both OFT and COFT exhibit consistent and stable convergence, with their synthesized face poses progressively aligning to the control landmarks. Conversely, ControlNet displays a sudden emergence of controllability at the 8-th epoch, corroborating the abrupt convergence previously reported in \cite{zhang2023adding}. This level of convergence stability seen in OFT offers practical advantages: training proceeds in a smoother and more predictable fashion, simplifying the search for effective hyperparameter settings.

\textbf{Quantitative comparison}. To assess the efficacy of controllable generation, we propose a metric based on control consistency. This metric involves extracting the control signal from the synthesized image and quantifying its similarity to the original control input. In the case of the C2I task, we report both IoU and F1 scores. For S2I, mean IoU as well as mean and overall accuracy are measured. For L2F, the average $\ell_2$ distance between the provided control landmarks and the landmarks predicted in the output is calculated. Detailed definitions of these metrics are provided in Appendix~\ref{app:settings}. All compared approaches are evaluated using their optimal hyperparameter configurations. According to the results in Table~\ref{table:control_gen}, OFT and COFT consistently achieve higher levels of control fidelity and precision in comparison to competing methods. Notably, adapter-based techniques such as T2I-Adapter and ControlNet display slower convergence and inferior ultimate performance. While LoRA surpasses ControlNet in the S2I scenario, it performs less favorably on C2I and L2F tasks. Broadly, our experiments indicate that LoRA’s maximum achievable performance remains limited, despite extensive hyperparameter optimization. By contrast, OFT continues to benefit from increased model capacity, with no clear plateau in observed performance. Our approach to quantitative evaluation represents a key contribution, enabling rigorous measurement of control efficacy in fine-tuned models, bounded only by the underlying segmentation or detection model’s accuracy.

\textbf{Qualitative comparison}. We conduct a qualitative assessment of OFT and COFT in comparison with existing state-of-the-art techniques, specifically ControlNet, T2I-Adapter, and LoRA. As illustrated by the sample outputs in Figure~\ref{fig:control_final}, both OFT and COFT consistently generate images that are not only realistic and high in fidelity but also allow for precise manipulation according to input controls. In the S2I task, LoRA is unable to generate images that correspond to the provided segmentation map, whereas ControlNet, OFT, and COFT demonstrate effective controllability. Notably, OFT and COFT deliver images with richer details and more accurate geometric configurations compared to ControlNet, despite utilizing significantly fewer model parameters. For the C2I task, OFT and COFT exhibit superior capability in generating semantically consistent images from rough Canny edge maps, outperforming both T2I-Adapter and LoRA. When addressing the L2F task, our approach offers the most reliable pose control for synthesized faces, even in cases of complex or challenging orientations. Across all three tasks, our results indicate that OFT and COFT surpass previous state-of-the-art approaches in image quality, underscoring the robustness of the OFT framework for controllable image generation. For a more thorough visual evaluation, additional qualitative results for each control task are presented in Appendix~\ref{app:control_exp}. Furthermore, Appendix~\ref{app:more_control_task} demonstrates OFT’s capability to excel in a broader range of control scenarios, such as dense pose to human body, sketch to image, and depth to image tasks.

\section{Concluding Remarks and Open Problems}

Inspired by the significant role that angular relationships among neurons play in defining visual semantics, we introduce a straightforward yet powerful fine-tuning approach—orthogonal finetuning—for exerting control over text-to-image diffusion models. This method is applied to two primary tasks: subject-driven generation and controllable generation. Relative to prior techniques, OFT achieves enhanced controllability and stable fine-tuning while requiring fewer trainable parameters. Crucially, OFT adds no extra computational burden during inference, making it suitable for practical deployment.

Nonetheless, OFT brings forth several compelling unresolved issues. One issue stems from its enforcement of orthogonality using Cayley parametrization, which necessitates a matrix inversion—posing scalability constraints. Although our adoption of block diagonal parametrization alleviates some of these concerns, efficiently computing matrix inverses within a differentiable framework remains unresolved. Another avenue arises from OFT's distinctive capacity for compositionality: the orthogonal matrices generated from separate OFT fine-tuning procedures can be multiplied to yield another orthogonal matrix. Investigating whether the knowledge from all downstream tasks is retained within this composite matrix constitutes an open research question. Lastly, OFT’s parameter efficiency is closely tied to its reliance on block diagonal structures, a strategy that introduces certain biases and restricts adaptability. Discovering methods to further enhance parameter efficiency with minimal bias and greater flexibility continues to be a vital and open challenge.